
\chapter{Proposed Adaptive Mesh Refinement Formulation}
\label{app:amr}

\section{Scope and status}
This appendix records the initial adaptive mesh refinement (AMR) design for
PSSFSS as of October 8, 2026. Electric-current and magnetic-current sheets are
both in scope. The objective is accurate generalized scattering matrices (GSMs)
and accurate scattering parameters of the assembled periodic structure, with
fewer unknowns than uniform refinement would require.  Note that AMR is probably
most useful for elements that employ unstructured meshes.  For structured meshes,
there are substantial efficiency gains obtained while filling the interaction
matrix due to exploiting the translational invariance of triangle-triangle
interactions.  These redundancies may be substantially reduced for meshes that
are auto-refined.  




\section{Source References}
References \cite{kim2017error,kim2018adaptive,kim2018evaluation,kim2022adaptive,pete:22}
motivate the design; periodic GSM aggregation and magnetic
duality are adaptations for PSSFSS rather than reported validation in those papers.



\subsection{Status terminology}
\begin{description}
\item[Implemented] Denotes existing reconstruction and tests.
\item[Agreed] Denotes a mathematical interpretation selected for initial development.
\item[Proposed] Denotes a candidate controller or integration strategy
\item[Open] Denotes a choice still requiring resolution.
\end{description}


\section{Reconstruction and Units}

\subsection{Modified RWG Normalization}
Let triangle $i$ have area $A_i$, centroid $\rc_i$, and vertex coordinates
$\r_{i1},\r_{i2},\r_{i3}$, all lengths expressed in meters.
For basis function $b$, let $\r_b^+$ and $\r_b^-$ be its free
vertices on its plus and minus triangles. For electric currents,
its electric surface-current contributions are
\begin{equation}
  \J_b^+(\r)  = \frac{I_b}{2A_b^+}(\r - \r_b^+),
  \qquad
 \J_b^-(\r)  = -\frac{I_b}{2A_b^-}(\r - \r_b^-).
 \label{eq:amr-rwg}
\end{equation}
Periodic half-bases additionally carry the phase factors described below.
There is no conventional RWG edge-length multiplier in
\eqref{eq:amr-rwg}: $I_b$ is the signed integrated current across the plus-side
edge, in \unit{A} (amperes). Current density has units \unit{A/m} and the impedance matrix has
units \unit{ohm} (ohms). For magnetic sources, replace $I_b$ by a coefficient $V_b$ in \unit{V} (volts)
and $\J$ by $\M$ in \unit{V/m}; the admittance matrix has units \unit{S} (siemens).
These expressions describe the effective basis used for reconstruction, including
the coefficient scaling in the main theory's RWG expansion.

\subsection{Continuity Equation}
For the existing $\exp(j\omega t)$ convention,
\begin{equation}
  q_{e,i} = \frac{j}{\omega}\nabla_s \cdot \mathbf J_i,
  \qquad
  q_{m,i} = \frac{j}{\omega}\nabla_s\cdot\mathbf M_i.
  \label{eq:amr-charge}
\end{equation}
The divergence is constant within a triangle. Thus each electric plus/minus
half-basis contributes $\pm j I_b/(\omega A_i)$, including its phase factor.
Electric charge density is in \unit{C/m^2}; magnetic charge density is in
\unit{V.s/m^2} (or \unit{Wb/m^2}). Both reconstructions use the same geometric
operations. 


\section{Centroid Interpretation of CWCD}
\subsection{Testing geometry}
The \emph{agreed initial interpretation} uses triangle-centroid current values,
not an exact line integral. For an interior edge shared by triangles $i$ and $n$,
let $\m_{in}$ be the edge midpoint and orient the testing path from $i$
to $n$. Define length-carrying testing vectors
\begin{equation}
  \dt_{i,in}= \m_{in} - \rc_i,
  \qquad
 \dt_{n,in} = \rc_n - \m_{in}.
 \label{eq:amr-test-vectors}
\end{equation}
Their identification with the dissertation's \cite{kim2017error}
testing vectors is a
\emph{proposed explicit geometric convention} to review before coding.
For the free vertex $\rfree_{i}$ opposite this edge,
$\dt_{i,in} = (\rc_i - \rfree_{i})/2$; similarly,
$\dt_{n,in} = (\rfree_{n} - \rc_n)/2$.
These identities connect the geometry to the half-vector factors in PSSFSS's
centroid testing expressions.

For excitation $a$, define
\begin{equation}
 B_{e,in}^{(a)} =
 \dt_{i,in} \cdot \J_i^{(a)}(\rc_i) + \dt_{n,in} \cdot \J_n^{(a)}(\rc_n),
 \qquad
 \Delta q_{e,in}^{(a)} = q_{e,i}^{(a)} - q_{e,n}^{(a)}.
 \label{eq:amr-centroid-test}
\end{equation}
The dot products are ordinary geometric contractions as defined for real vectors; they do
not conjugate the complex current. Define $B_m$ and $\Delta q_m$ analogously for magnetic quantities.
Reversing the edge orientation changes both signs but not their magnitudes.
RWG currents are affine within each triangle: centroid evaluation generally
differs from exact integration along each half-path. This discrete interpretation
must not be silently replaced by path integration.

\subsection{Electric and Magnetic Interface Weights}
The \emph{agreed interface policy} uses the material factors appearing in
PSSFSS's source-plane potential Green functions. Let $\mu_s,\epsilon_s$ and
$\mu_{s+1}, \epsilon_{s+1}$ be the absolute, possibly complex permeabilities and permittivities
immediately to the left and right of the sheet. Define
\begin{equation}
  \tilde{\mu} = \frac{2 \mu_s \mu_{s+1}}{\mu_s + \mu_{s+1}},
  \qquad
  \bar{\epsilon} = \frac{\epsilon_s + \epsilon_{s+1}}{2}.
 \label{eq:amr-interface-materials}
\end{equation}
These are the harmonic mean of permeability and arithmetic mean of permittivity.
For electric sources, Chapter~\ref{chap:gfstratified}'s accelerated Green functions factor out
$\tilde{\mu}$ from the vector potential and $1/\bar{\epsilon}$ from
the scalar potential (see \eqref{eq:finalEGF}).
Their local, large-spatial-frequency behavior motivates the same factors in
CWCD; this is an estimator adaptation, not a replacement of the full layered
Green functions.

For positive lossless parameters, one may define
\begin{equation}
  k_{\text{eff}} = \omega\sqrt{\tilde{\mu}\,\bar{\epsilon}},\qquad
  \eta_{\text{eff}}= \sqrt{\frac{\tilde{\mu}}{\bar{\epsilon}}},\qquad
  k_{\text{eff}} \eta_{\text{eff}} = \omega\tilde{\mu},\quad
 \frac{k_{\text{eff}}}{\eta_{\text{eff}}} = \omega \bar{\epsilon}.
 \label{eq:amr-effective-k-eta}
\end{equation}
There is no need to compute either square root for the indicator. Use directly
\begin{align}
  s_{e,in}^{(a)} &=
                   \left|\omega\tilde{\mu}\,B_{e,in}^{(a)}\right|
                   + \left|\frac{\Delta q_{e,in}^{(a)}}{\bar{\epsilon}}\right|,
                   \label{eq:amr-electric-score} \\
  s_{m,in}^{(a)} &=
                   \left|\omega\bar{\epsilon}\,B_{m,in}^{(a)}\right|
                   +\left|\frac{\Delta q_{m,in}^{(a)}}{\tilde{\mu}}\right|.
                   \label{eq:amr-magnetic-score}
\end{align}
The electric score defined above has units \unit{V/m}; the magnetic score has units \unit{A/m}.
CWCD adds a current term to a charge jump, rather than multiplying the jump by
current magnitude. Material coefficients remain inside each absolute value,
so the scores are nonnegative for complex material parameters and no
square-root branch selection is required.

\subsection{Why magnetic permeability is harmonically averaged}
The magnetic formulation adds independent left-looking and right-looking
potential Green functions, as described in Chapter~\ref{chap:mom}'s magnetic-source
formulation (see \eqref{eq:GFxxGpsileftright}). Each side sees
its own one-sided structure. The leading current weights are proportional to
$\omega\epsilon_s$ and $\omega\epsilon_{s+1}$, whereas the charge weights are
proportional to $1/\mu_s$ and $1/\mu_{s+1}$. Consequently,
\begin{equation}
 \frac12 \left(\frac{1}{\mu_s}+\frac{1}{\mu_{s+1}}\right)
  = \frac{1}{\tilde{\mu}},
 \label{eq:amr-reciprocal-mu}
\end{equation}
rather than $1/\bar{\mu}$ with $\bar{\mu}=(\mu_s+\mu_{s+1})/2$. The accelerated
magnetic Green functions in Chapter~\ref{chap:gfstratified} likewise contain $\bar{\epsilon}$ and
$1/\tilde{\mu}$. Using averages rather than sums omits a common factor
of two from the magnetic weights; this convention must be applied consistently
to any solution-based normalization. An incident-field normalization would
require an explicit choice of that overall scale.

For equal materials, the scores reduce to the homogeneous electric and magnetic
forms. For equal permeabilities, arithmetic and harmonic means of permeability
coincide, including the common nonmagnetic dielectric case. Remote layers still
affect the solved currents and charges, but are not included in these local
material weights. The selected weights do not establish CWCD accuracy for every
material contrast or loss level; validate them by mesh convergence.

In implementation, \texttt{FillZY.jl} uses relative versions of these averages:
multiply the harmonic relative permeability by $\mu_0$ and the arithmetic
relative permittivity by $\epsilon_0$ before applying SI reconstruction outputs.
Zero or singular effective parameters require explicit handling; do not silently
regularize them or substitute magnitudes before averaging.

For a triangle with three defined neighbour contributions, let
\begin{equation}
 E_i^{(a)}=\frac{1}{3}\sum_{m=1}^3 s_{i,n(i,m)}^{(a)},\qquad
 \eta_i^{(a)}=\frac{E_i^{(a)}}{D^{(a)}}.
 \label{eq:amr-local-indicator}
\end{equation}
A physical boundary cannot be inserted into this three-neighbour
sum without a specified boundary convention.
In \eqref{eq:amr-local-indicator} $D^{(a)}>0$ is a normalization constant with the same units as $E$.
There are several possibilities for its definition. The dissertation \cite[Eq.~(5.7)]{kim2017error}
normalizes by incident field magnitude; the published CWCD equation \cite[Eq.~(5)]{kim2018evaluation}
uses a solution-based scale. We will follow \cite{kim2018evaluation} and choose
\begin{equation}
  \label{eq:Ddef}
  D^{(a)} = \max_i E_i^{(a)}.
\end{equation}
In the remote possibility that $D^{(a)} = 0$ an error condition will be raised.


\section{Periodic Neighbours and Physical Boundaries}
For a translated periodic image displaced by
$\R_{pq} = p \s_1 + q \s_2$, the existing convention is
\begin{equation}
 \J(\r + \R_{pq}) =  e^{-j(p\psi_1+q\psi_2)} \J(\r).
 \label{eq:amr-periodic-phase}
\end{equation}
The same phase relation applies to charge and magnetic current. Translate the
neighbor centroid and free vertex into the image adjacent to the chosen edge,
and apply the associated phase to its reconstructed field quantities before
forming \eqref{eq:amr-centroid-test}. Reversing the image translation reverses
the phase exponent. Do not apply the reconstruction half-basis factor twice:
reconstruction produces fields in the stored cell; image comparison is a
separate operation. Verify this convention against RWG edge codes and tests.

A physical metal or aperture boundary has no conducting or aperture neighbour
across it. Its score treatment is \emph{open}; excluding the missing edge,
constructing an exterior continuation, and using a boundary-specific score are
not equivalent. A periodic seam is not a physical edge. Refinement must preserve
paired boundary segmentation and prevent hanging nodes.

\section{Excitations and marking}
The \emph{proposed initial GSM objective} uses all retained incident TE and TM
Floquet modes from both sides of each Gblock, including evanescent modes.
The local system's factorization is shared across these right-hand sides.
Separately normalized indicators can be accumulated without storing every
solution:
\begin{equation}
 \eta_i^{\text{GSM}}=\max_{a\in\mathcal E}\eta_i^{(a)}.
 \label{eq:amr-mode-maximum}
\end{equation}
This is very simple to implement, but gives weak and strongly excited responses comparable influence,
which may not be the best choice.  Instead, a better choice might be to base AMR on the actual total current induced via a given external illumination.
This would be more complicated to implement and is postponed for later examination.
As a brief summary, this latter method would reconstruct the actual current
$\J^{(b)}=\sum_a u_a^{(b)} \J^{(a)}$ for an external illumination
$b$, using incident modal amplitudes recovered from cascading. Compute CWCD
only once \emph{after} this coherent sum.
This is an optional alternative objective, not a requirement for
the initial GSM implementation.

A proposed area-weighted marking score is
\begin{equation}
 M_i = A_i\eta_i^{\text{GSM}}, \qquad  \mathcal F = \{i: M_i> \zeta\max_j M_j\}, \quad 0 < \zeta < 1.
 \label{eq:amr-marking}
\end{equation}
The papers use $\zeta=0.7$ in particular experiments; this is not a universal
accuracy tolerance. Area weighting is a controller choice, distinct from CWCD.
It reduces the dominance of shrinking cells near physical edge singularities.
A zero maximum should produce no marks. A decreasing maximum charge jump is not
a necessary condition for improved scattering accuracy near singular edges.

The candidate refinement footprint adds edge midpoints to marked cells and
centroids to neighbouring cells sharing a vertex, followed by mesh smoothing via
constrained Laplace iteration, then constrained triangulation.  Periodic partners must be
included consistently. Boundary geometry, holes, conformity, and triangle quality
must be preserved; any smoothing must respect physical and periodic constraints.
The specific mesher and footprint remain \emph{proposed}. 

The adaptation sample selection and treatment of a changing mesh during a fast
interpolated sweep remain \emph{open}.  However, the initial proposed sampling
plan is the following heuristic: 
Perform mesh refinement at only the highest requested analysis frequency, and
at the most extreme requested scan angle (or incremental phasing, if phases
rather than steering angles are specified by the user).  The most ``extreme''
steering angle will mean that with maximum $|\theta|$  (if not unique, then
select one of these with maximum $|\phi|$). For incremental phasings,
the most ``extreme'' phasing will be any one maximizing $|\psi_1|^2 + |\psi_2|^2$.
Eventually, the plan is to allow the user to specify one or more
(scan condition, frequency) combinations at which to perform adaptive meshing,
obtaining a single, common mesh that is then used for all requested sweeps.
If none are specified, then the heuristic described above will be used.

\section{Validation, Stopping Criteria}
Use analytic reconstruction tests and periodic image tests before integration.
Compare adaptive results against converged uniformly refined references at
representative frequencies, scan conditions, and both source classes. Examine
complex GSM entries and final scattering parameters; indicator reduction alone
is not evidence of the required accuracy. Modal truncation error must be
controlled separately from mesh discretization error.

Stopping criteria will include
\begin{itemize}
\item S-parameter convergence,
\item Maximum allowed number of triangles exceeded.
\end{itemize}
The AMR algorithm will terminate whenever either of these is criteria is met.
Exact thresholds and comparison
norms remain \emph{open}. A resource-limited stop must not be reported as
convergence. 

